\chapter{GPU Kernels with Triton}\label{app:triton}

Throughout the book we have argued with rooflines: decode is memory-bound, softmax is
memory-bound, attention without tiling moves $O(T^2)$ bytes, and a weight-only int4 GEMM
is fast only if dequantization happens \emph{inside} the kernel. This appendix makes those
arguments executable. We write the kernels ourselves in Triton (Tillet et al., 2019), a
Python-embedded language that compiles block-level programs to GPU code. Triton is what
\code{torch.compile}'s Inductor backend emits for most fused operations, so reading
Triton is also how you inspect what the compiler did to your model.

Every kernel here has a PyTorch reference in \ident{llmsys/kernels/reference.py}. The
Triton versions live in \ident{llmsys/kernels/triton_kernels.py}. Triton ships with an
\emph{interpreter}: set \code{TRITON\_INTERPRET=1} before importing \code{triton}, and
kernels run on the CPU with NumPy semantics. It is slow, but it means correctness can be
tested on a laptop, and it is how this appendix's lab is checked.

\begin{figure}[t]
\centering
\begin{tikzpicture}[
  stage/.style={draw=sysblue,fill=sysbluebg,rounded corners=2pt,minimum width=2.3cm,
    minimum height=1.15cm,align=center,font=\small},
  num/.style={circle,fill=sysblue,text=white,font=\bfseries\small,inner sep=1.5pt},
  arr/.style={-Latex,thick,sysblue!70}]
\node[stage] (s1) {Programming\\model};
\node[stage,right=0.35cm of s1] (s2) {Vector add\\\& debugging};
\node[stage,right=0.35cm of s2] (s3) {Fused softmax\\\& RMSNorm};
\node[stage,right=0.35cm of s3] (s4) {Tiled\\matmul};
\node[stage,below=0.8cm of s2] (s5) {W8A16\\GEMM};
\node[stage,right=0.35cm of s5] (s6) {FlashAttention\\forward};
\node[stage,right=0.35cm of s6] (s7) {Profiling \&\\integration};
\foreach \i in {1,...,7}{\node[num] at (s\i.north west) {\i};}
\draw[arr] (s1)--(s2); \draw[arr] (s2)--(s3); \draw[arr] (s3)--(s4);
\draw[arr] (s4.south) |- ++(0,-0.4) -| (s5.north);
\draw[arr] (s5)--(s6); \draw[arr] (s6)--(s7);
\end{tikzpicture}
\caption{Roadmap of this appendix: from the programming model to kernels that implement
\cref{ch:attention}'s FlashAttention and \cref{ch:efficient}'s quantized GEMM.}
\label{fig:triton:roadmap}
\end{figure}

% =====================================================================
\section{Why write a kernel at all?}
A PyTorch eager program launches one kernel per operator. Each kernel reads its inputs
from HBM and writes its outputs back to HBM. For matmuls this is fine, since their
arithmetic intensity is high. For chains of elementwise and reduction operations
(softmax, norms, activation functions, dequantization, rotary embeddings) the
intermediate tensors make round trips through memory, and the chain runs at a small
fraction of what the hardware could do. A \emph{fused} kernel loads each input once,
keeps intermediates in registers or shared memory, and writes the result once.

\begin{keyidea}
Write a custom kernel when (1) a sequence of memory-bound operations can be fused, (2) a
computation's working set can be \emph{tiled} into on-chip memory to change its
asymptotic memory traffic (FlashAttention), or (3) a data format is not supported by
library kernels (packed int4 weights, custom KV layouts). Otherwise, cuBLAS, cuDNN and
\code{torch.compile} usually beat a hand-written kernel.
\end{keyidea}

% =====================================================================
\section{The Triton programming model}

\subsection{From CUDA threads to Triton programs}
CUDA exposes a three-level hierarchy. A \emph{grid} of thread blocks; each block is
scheduled onto one streaming multiprocessor (SM) and contains up to 1024 threads grouped
into 32-wide warps; each thread owns registers, and threads in a block share up to a few
hundred KB of shared memory (SMEM). The programmer reasons about individual threads:
which element each one loads, how they cooperate through SMEM, and how to synchronise.

Triton removes the thread level. You write a \emph{program} that operates on whole
\emph{blocks} (tiles) of data as if they were small arrays, and launch a grid of program
instances. The compiler decides how to spread each block across threads and warps, when
to stage data through shared memory, how to vectorise memory accesses, and how to issue
tensor-core instructions for \code{tl.dot}. \Cref{fig:triton:hierarchy} contrasts the
two views.

\begin{figure}[t]
\centering
\begin{tikzpicture}[font=\scriptsize,
  g/.style={draw,rounded corners=2pt,fill=#1,align=center},
  arr/.style={-Latex,thick}]
% CUDA
\node[font=\small\bfseries] at (2.4,3.3) {CUDA};
\node[g=gray!10,minimum width=4.8cm,minimum height=2.9cm] (grid) at (2.4,1.4) {};
\node[anchor=north west] at (grid.north west) {grid};
\foreach \i/\x in {0/0.9,1/2.4,2/3.9}{
  \node[g=sysblue!15,minimum width=1.3cm,minimum height=2.1cm] (b\i) at (\x,1.2) {};
  \node[anchor=north] at (b\i.north) {block \i};
  \foreach \w in {0,1,2}{\node[g=keygreen!25,minimum width=1.05cm,minimum height=0.36cm,font=\tiny] at (\x,1.25-\w*0.48) {warp (32)};}}
\node[below=0.05cm of grid,align=center,text width=4.6cm] {you index \emph{threads}; manage SMEM, syncs, MMA fragments};
% Triton
\node[font=\small\bfseries] at (8.6,3.3) {Triton};
\node[g=gray!10,minimum width=4.8cm,minimum height=2.9cm] (tgrid) at (8.6,1.4) {};
\node[anchor=north west] at (tgrid.north west) {grid of programs};
\foreach \i/\x in {0/7.1,1/8.6,2/10.1}{
  \node[g=sysblue!15,minimum width=1.3cm,minimum height=2.1cm] (p\i) at (\x,1.2) {};
  \node[anchor=north] at (p\i.north) {pid \i};
  \node[g=napkin!25,minimum width=1.05cm,minimum height=1.3cm,align=center] at (\x,0.95) {block\\of data\\(tile)};}
\node[below=0.05cm of tgrid,align=center,text width=4.6cm] {you index \emph{blocks}; compiler maps to warps, SMEM, tensor cores};
\draw[arr,labpurple] (grid.east) -- (tgrid.west) node[midway,above] {abstract};
\end{tikzpicture}
\caption{CUDA's thread-level view versus Triton's block-level view. A Triton program
corresponds to a CUDA thread block. Inside it you write NumPy-like operations on tiles,
with \code{num\_warps} controlling how many warps execute it.}
\label{fig:triton:hierarchy}
\end{figure}

\subsection{The memory hierarchy a kernel sees}
Kernel design is mostly about \emph{where data lives}. For an H100 SXM:
\begin{center}\small
\begin{tabularx}{\linewidth}{l>{\raggedright\arraybackslash}Xll}
\toprule
level & capacity & scope & managed by \\
\midrule
registers & 256\,KB per SM (64k $\times$ 32-bit) & one thread & compiler \\
SMEM / L1 & 256\,KB per SM combined, up to 228\,KB usable as SMEM & thread block & you / compiler \\
L2 cache & 50\,MB & whole GPU & hardware \\
HBM3 & 80\,GB at 3.35\,TB/s & whole GPU & you \\
\bottomrule
\end{tabularx}
\end{center}
There are 132 SMs. Each step down the hierarchy is roughly an order of magnitude more
capacity and an order of magnitude less bandwidth per byte. Tiles must therefore be
sized so that their working set (inputs, accumulators, pipeline buffers) fits in
registers and SMEM, while still leaving enough room for several programs per SM to hide
latency (\emph{occupancy}). In Triton you control this through block sizes,
\code{num\_warps} and \code{num\_stages}. Exceed the register budget and the compiler
spills to local memory, which lives in HBM, and performance collapses
(\cref{fig:triton:memory}).

\begin{figure}[t]
\centering
\begin{tikzpicture}[font=\scriptsize]
\foreach \y/\w/\c/\t/\r in {0/2.0/fcompute!45/registers/fastest; private,
                            -0.8/3.6/fcompute!25/shared memory (SMEM)/explicit tiles,
                            -1.6/5.2/fmemory!30/L2 cache (50 MB)/implicit reuse across SMs,
                            -2.4/6.8/fmemory!55/HBM (80 GB{,} 3.35 TB/s)/the roofline's denominator}{
  \node[draw,fill=\c,minimum width=\w cm,minimum height=0.65cm] at (0,\y) {\t};
  \node[anchor=west,text=srcgray] at (3.7,\y) {\r};}
\draw[-Latex,thick] (-4.0,-2.6) -- (-4.0,0.2) node[midway,left,align=right] {faster,\\smaller};
\draw[-Latex,thick,warnred,dashed] (-1.2,0) to[bend right=40] node[left,align=right,pos=0.6] {spills} (-2.8,-2.4);
\end{tikzpicture}
\caption{The memory hierarchy as seen by one kernel. Triton tiles live in registers and
SMEM. Register spills silently fall all the way to HBM.}
\label{fig:triton:memory}
\end{figure}

\subsection{A first kernel: vector addition}
\Cref{lst:triton:add} adds two vectors. The \code{@triton.jit} kernel runs once per
program instance:
\begin{itemize}
  \item \code{tl.program\_id(0)} tells the instance which block it owns;
  \item \code{tl.arange(0, BLOCK)} creates a vector of offsets within the block;
  \item \code{mask} guards the ragged last block (\cref{fig:triton:add});
  \item pointers are integers plus offsets, so \code{x\_ptr + offsets} is a block of
        addresses and \code{tl.load} gathers them.
\end{itemize}
\code{BLOCK} is a \code{tl.constexpr}: it is known at compile time (and must be a power
of two), so the compiler can specialise register allocation and vector widths. A new
value triggers a recompilation, which is cached.

\begin{lstlisting}[caption={Vector addition in Triton, with its launcher.},label={lst:triton:add}]
@triton.jit
def _add_kernel(x_ptr, y_ptr, out_ptr, n_elements, BLOCK: tl.constexpr):
    pid = tl.program_id(axis=0)                  # which block am I?
    offsets = pid * BLOCK + tl.arange(0, BLOCK)  # BLOCK element indices
    mask = offsets < n_elements                  # guard the ragged tail
    x = tl.load(x_ptr + offsets, mask=mask)
    y = tl.load(y_ptr + offsets, mask=mask)
    tl.store(out_ptr + offsets, x + y, mask=mask)

def vector_add(x, y, block=1024):
    out = torch.empty_like(x)
    n = x.numel()
    grid = (triton.cdiv(n, block),)              # number of program instances
    _add_kernel[grid](x, y, out, n, BLOCK=block)
    return out
\end{lstlisting}

\begin{figure}[t]
\centering
\begin{tikzpicture}[x=0.32cm,y=0.5cm,font=\scriptsize]
\foreach \i in {0,...,37}{
  \pgfmathtruncatemacro{\p}{\i/10}
  \ifnum\p=0 \def\c{sysblue!30}\fi
  \ifnum\p=1 \def\c{keygreen!35}\fi
  \ifnum\p=2 \def\c{napkin!35}\fi
  \ifnum\p=3 \def\c{labpurple!30}\fi
  \fill[\c] (\i,0) rectangle (\i+1,1);}
\foreach \i in {38,39}{\fill[pattern=north east lines,pattern color=warnred] (\i,0) rectangle (\i+1,1);}
\draw[step=1] (0,0) grid (40,1);
\foreach \p in {0,1,2,3}{\draw[decorate,decoration={brace,amplitude=3pt}] (\p*10,1.2) -- (\p*10+10,1.2)
  node[midway,above=3pt] {\texttt{pid=\p}};}
\node[below,align=center,text=warnred] at (39,-0.1) {masked\\($\ge n$)};
\node[anchor=west] at (0,-1.0) {$n=38$, \texttt{BLOCK}=10, grid $=\lceil 38/10\rceil=4$ programs};
\end{tikzpicture}
\caption{Each program instance handles one contiguous block of \code{BLOCK} elements.
The last block is partially out of bounds and is masked: masked loads return
\code{other} (default undefined, so always pass it when the value matters), and masked
stores are dropped.}
\label{fig:triton:add}
\end{figure}

\subsection{Debugging with the interpreter}
Kernels compile lazily on first launch, with arguments' dtypes and \code{constexpr}
values as the cache key. Errors at that point tend to be cryptic. Two habits help.
First, develop with \code{TRITON\_INTERPRET=1}: kernels then execute eagerly in Python,
so you can \code{print} a tile or drop into a debugger inside the kernel, on a machine
without a GPU. Second, test against a PyTorch reference on awkward shapes (prime sizes,
sizes smaller than one block, non-contiguous strides) before measuring speed.

\begin{lstlisting}[style=shell]
TRITON_INTERPRET=1 python -m pytest -q labs/labA_triton     # correctness, CPU
python labs/labA_triton/bench.py --hw H100-SXM               # speed, GPU
\end{lstlisting}

\begin{pitfall}[The interpreter flag must be set first]
\code{TRITON\_INTERPRET} is read when \code{@triton.jit} decorates a function, so it must
be in the environment \emph{before} the module defining your kernels is imported. Our
lab's test file sets it at the top when no GPU is present.
\end{pitfall}

% =====================================================================
\section{Fused softmax and RMSNorm}\label{sec:triton:softmax}
A row-wise softmax in eager PyTorch, written out (\code{x - max}, \code{exp},
\code{sum}, divide), launches several kernels, each streaming the whole matrix through
HBM. \Cref{lst:triton:softmax} assigns one program per row, loads the row once,
performs every step in registers, and writes the result once.

\begin{lstlisting}[caption={Fused row softmax: one program per row.},label={lst:triton:softmax}]
@triton.jit
def _softmax_kernel(x_ptr, out_ptr, n_cols, x_row_stride, out_row_stride,
                    BLOCK: tl.constexpr):
    row = tl.program_id(0)
    cols = tl.arange(0, BLOCK)
    mask = cols < n_cols
    x = tl.load(x_ptr + row * x_row_stride + cols, mask=mask, other=-float("inf"))
    x = x.to(tl.float32)
    z = x - tl.max(x, axis=0)             # padding is -inf, so it never wins the max
    num = tl.exp(z)                       # exp(-inf) = 0: padding adds nothing
    den = tl.sum(num, axis=0)
    tl.store(out_ptr + row * out_row_stride + cols, num / den, mask=mask)
\end{lstlisting}

Note how the padding value is chosen so that masked lanes are neutral for both
reductions: $-\infty$ for the max, and therefore $e^{-\infty}=0$ for the sum.

\begin{figure}[t]
\centering
\begin{tikzpicture}[font=\scriptsize,
  k/.style={draw,fill=fcompute!25,rounded corners=1pt,minimum width=1.25cm,minimum height=0.55cm},
  h/.style={draw,fill=fmemory!30,minimum width=10.4cm,minimum height=0.4cm}]
\node[anchor=east,font=\small] at (-0.2,1.6) {eager};
\foreach \n/\x in {max/0.65,sub/2.15,exp/3.65,sum/5.15,div/6.65}{\node[k] (e\n) at (\x,1.6) {\n};}
\node[h] (hbm) at (5.2,0.3) {HBM};
\foreach \n in {max,sub,exp,sum,div}{\draw[-Latex,warnred] (e\n.south) ++(-0.2,0) -- ++(0,-0.82);
  \draw[Latex-,keygreen!60!black] (e\n.south) ++(0.2,0) -- ++(0,-0.82);}
\node[anchor=west,align=left] at (7.6,1.6) {5 kernels, $\approx$5 reads + 4 writes\\of the matrix (sum writes one vector)};
\node[anchor=east,font=\small] at (-0.2,-1.1) {fused};
\node[k,minimum width=3.0cm] (f) at (1.6,-1.1) {load $\to$ max, exp, sum, div};
\draw[Latex-,warnred] (f.north) ++(-0.4,0) -- ++(0,0.75);
\draw[-Latex,keygreen!60!black] (f.north) ++(0.4,0) -- ++(0,0.75);
\node[anchor=west,align=left] at (3.4,-1.1) {1 kernel, 1 read + 1 write; intermediates stay in registers};
\end{tikzpicture}
\caption{Memory traffic of an unfused versus a fused softmax. Red arrows are reads from
HBM, green arrows are writes. For a memory-bound operation, run time is proportional to
the arrows.}
\label{fig:triton:softmax}
\end{figure}

\begin{napkinmath}[Softmax over a 4096 $\times$ 8192 fp32 matrix on an H100]
The matrix is 134\,MB. Fused: one read plus one write, 268\,MB, so
$\ge 80\,\mu$s at 3.35\,TB/s. The five-kernel eager version moves roughly
$9\times134$\,MB\,$\approx1.2$\,GB, so $\ge360\,\mu$s. PyTorch's \code{torch.softmax} is
itself a fused kernel; our Triton kernel should match it, and the lab's benchmark
reports achieved GB/s against the 3.35\,TB/s peak. A good fused kernel reaches 80--90\%.
\end{napkinmath}

The one-row-per-program design assumes a row fits in registers. Beyond roughly 32k
fp32 columns, register pressure forces spills, and the kernel must loop over chunks of
the row. It then either does two passes (max, then sum) or keeps \emph{running}
statistics with the online-softmax rescaling of \cref{sec:triton:flash}. RMSNorm
(\cref{ch:modern}) has the same shape: one reduction (the mean of squares) followed by
an elementwise scale. The lab implements it exactly like softmax.

\subsection{A second fusion: the SwiGLU activation}
Qwen3's MLP computes $\mathrm{down}\big(\SiLU(\mathrm{gate}(x))\odot\mathrm{up}(x)\big)$
(\cref{ch:modern}). The two projections are GEMMs, but between them sit three
memory-bound steps: the SiLU, the elementwise product, and the dtype handling. Fusing
them into one pass is the simplest useful kernel there is (\cref{lst:triton:swiglu}).

\begin{lstlisting}[caption={Fused SiLU-and-multiply (\ident{llmsys/kernels/triton_kernels.py}).},
label={lst:triton:swiglu}]
@triton.jit
def _swiglu_kernel(g_ptr, u_ptr, out_ptr, n_elements, BLOCK: tl.constexpr):
    offs = tl.program_id(0) * BLOCK + tl.arange(0, BLOCK)
    mask = offs < n_elements
    g = tl.load(g_ptr + offs, mask=mask, other=0.0).to(tl.float32)
    u = tl.load(u_ptr + offs, mask=mask, other=0.0).to(tl.float32)
    silu = g / (1.0 + tl.exp(-g))                 # g * sigmoid(g)
    tl.store(out_ptr + offs, silu * u, mask=mask)
\end{lstlisting}

\begin{napkinmath}[SwiGLU fusion during prefill]
Take Qwen3-8B ($d_{\text{ff}}=12288$) and a prefill chunk of 8192 tokens in bf16. Each of
\code{gate} and \code{up} is $8192\cdot12288\cdot2$\,B$=201$\,MB. Unfused (\code{silu},
then \code{mul}) reads 201\,MB and writes 201\,MB, then reads 402\,MB and writes
201\,MB: about 1.0\,GB, or 300\,$\mu$s at 3.35\,TB/s, per layer. Fused, it reads 402\,MB
and writes 201\,MB: 603\,MB, or 180\,$\mu$s. Over 36 layers that saves about 4\,ms per
chunk. An even better option is to fuse the activation into the \emph{epilogue} of the
gate/up GEMM, so the 201\,MB intermediates never reach HBM at all.
\end{napkinmath}

\subsection{A backward kernel: RMSNorm}
Training needs gradients, and fused forward kernels need fused backward kernels.
For one row, let $r=(\frac1n\sum_j x_j^2+\epsilon)^{-1/2}$, $\hat x=r\,x$ and
$y=w\odot\hat x$. Given the upstream gradient $\vect{dy}$, define $\vect g=\vect{dy}\odot w$.
Then
\begin{equation}
  \frac{\partial \mathcal L}{\partial x} \;=\; r\Big(\vect g - \hat x\,\frac{1}{n}\sum_j g_j\hat x_j\Big),
  \qquad
  \frac{\partial \mathcal L}{\partial w} \;=\; \sum_{\text{rows}} \vect{dy}\odot\hat x .
  \label{eq:triton:rmsbwd}
\end{equation}
The input gradient is again a one-row reduction followed by elementwise work: one
program per row, exactly like the forward. The weight gradient is a reduction
\emph{across} rows. The kernel writes per-row partial products and a second, tiny
reduction sums them. Production kernels write one partial per program (not per row)
and use a second Triton kernel or atomics. Note that $r$ is \emph{recomputed} from $x$
rather than saved: recomputing a reduction over data we must load anyway costs nothing,
while storing it costs memory.

\begin{lstlisting}[caption={Fused RMSNorm backward and its autograd wrapper.},
label={lst:triton:rmsbwd}]
@triton.jit
def _rmsnorm_bwd_kernel(x_ptr, w_ptr, dy_ptr, dx_ptr, gw_ptr, n_cols, stride, eps,
                        BLOCK: tl.constexpr):
    row = tl.program_id(0)
    cols = tl.arange(0, BLOCK)
    mask = cols < n_cols
    x = tl.load(x_ptr + row * stride + cols, mask=mask, other=0.0).to(tl.float32)
    w = tl.load(w_ptr + cols, mask=mask, other=0.0).to(tl.float32)
    dy = tl.load(dy_ptr + row * stride + cols, mask=mask, other=0.0).to(tl.float32)
    rstd = 1.0 / tl.sqrt(tl.sum(x * x, axis=0) / n_cols + eps)   # recomputed
    xhat = x * rstd
    g = dy * w
    dx = rstd * (g - xhat * (tl.sum(g * xhat, axis=0) / n_cols))
    tl.store(dx_ptr + row * stride + cols, dx, mask=mask)
    tl.store(gw_ptr + row * stride + cols, dy * xhat, mask=mask)  # dW partials

class TritonRMSNorm(torch.autograd.Function):
    @staticmethod
    def forward(ctx, x, weight, eps=1e-6):
        ctx.save_for_backward(x, weight); ctx.eps = eps
        return rmsnorm(x, weight, eps)
    @staticmethod
    def backward(ctx, dy):
        x, weight = ctx.saved_tensors
        dx, dw = rmsnorm_backward(x, weight, dy, ctx.eps)
        return dx, dw, None
\end{lstlisting}
The lab's tests check both gradients against PyTorch autograd, under the interpreter.

\begin{systemslens}[Fusion is where \code{torch.compile} earns its keep]
Inductor, \code{torch.compile}'s default backend, performs exactly this kind of
fusion automatically. It groups pointwise and reduction nodes of the captured graph and
emits Triton kernels like \cref{lst:triton:softmax}. Run with
\code{TORCH\_LOGS=output\_code} to see the generated Triton for your model. Hand-written
kernels still win where the compiler cannot see the opportunity, such as tiling
attention, custom quantized formats, or cross-operation algebraic rewrites.
\end{systemslens}

% =====================================================================
\section{Tiled matrix multiplication}\label{sec:triton:matmul}
\subsection{Tiling and arithmetic intensity}
Computing $C=AB$ with $A\in\R^{M\times K}$, $B\in\R^{K\times N}$, each program produces
one $B_M\times B_N$ tile of $C$. It loops over $K$ in slabs of $B_K$, loading a
$B_M\times B_K$ tile of $A$ and a $B_K\times B_N$ tile of $B$ and accumulating their
product in fp32 registers (\cref{fig:triton:tiling}). Per slab it moves
$(B_MB_K+B_KB_N)$ elements and performs $2B_MB_NB_K$ FLOPs, so its intensity with
respect to the loads is
\begin{equation}
  I_{\text{tile}} = \frac{2B_MB_NB_K}{(B_M+B_N)B_K\cdot b} = \frac{2B_MB_N}{(B_M+B_N)\,b},
\end{equation}
where $b$ is bytes per element. With $128\times128$ bf16 tiles,
$I=2\cdot128^2/(256\cdot2)=64$ FLOP/byte. That is below the H100's ridge of
$\approx295$ \emph{if those loads went to HBM}. They do not all go there: tiles of $A$
are reused by every program in the same row of the output grid, and tiles of $B$ by
every program in the same column. The L2 cache serves most re-reads, which is why
program \emph{ordering} matters (next subsection). Larger tiles raise $I$ but consume
more registers and SMEM, lowering occupancy. Autotuning searches this trade-off.

\begin{figure}[t]
\centering
\begin{tikzpicture}[x=0.38cm,y=0.38cm,font=\scriptsize]
% A
\draw (0,0) rectangle (8,-6); \node at (4,0.6) {$A$ ($M\times K$)};
\fill[sysblue!35] (0,-2) rectangle (8,-4);
\fill[sysblue!75] (2,-2) rectangle (4,-4);
\node[text=white] at (3,-3) {$B_K$};
% B
\draw (10,0) rectangle (16,-8); \node at (13,0.6) {$B$ ($K\times N$)};
\fill[keygreen!35] (12,0) rectangle (14,-8);
\fill[keygreen!75] (12,-2) rectangle (14,-4);
% C
\draw (18,0) rectangle (24,-6); \node at (21,0.6) {$C$ ($M\times N$)};
\fill[napkin!70] (20,-2) rectangle (22,-4);
\node[align=center] at (21,-7.2) {program $(p_m,p_n)$ owns\\one $B_M\times B_N$ tile};
\draw[-Latex,thick] (2.2,-6.5) -- (7.8,-6.5) node[midway,below,align=center]{loop over $K$ in $B_K$ slabs};
\node at (9.1,-3) {$\times$}; \node at (17,-3) {$=$};
\draw[decorate,decoration={brace,amplitude=3pt,mirror}] (-0.2,-2) -- (-0.2,-4) node[midway,left=3pt] {$B_M$};
\draw[decorate,decoration={brace,amplitude=3pt,mirror}] (12,-8.15) -- (14,-8.15) node[midway,below=3pt] {$B_N$};
\end{tikzpicture}
\caption{Tiled GEMM. The program owning the orange tile of $C$ streams the light-blue
row panel of $A$ and the light-green column panel of $B$ through on-chip memory, one
$B_K$-wide slab at a time (dark), accumulating in fp32 registers.}
\label{fig:triton:tiling}
\end{figure}

\subsection{Program ordering for L2 reuse}
Programs are launched in increasing \code{pid} order, and only as many run concurrently
as fit on the SMs. With naive row-major ordering, the concurrently running programs
cover one full row of output tiles. They share one panel of $A$, but each needs a
\emph{different} panel of $B$, so the working set is the whole of $B$. \emph{Grouped}
(``swizzled'') ordering walks $G$ rows of tiles in column-major order within each
group, so concurrent programs share a few panels of both $A$ and $B$
(\cref{fig:triton:swizzle}). The remapping is pure integer arithmetic on the program id,
at the top of \cref{lst:triton:matmul}.

\begin{figure}[t]
\centering
\begin{tikzpicture}[x=0.5cm,y=0.5cm,font=\scriptsize]
\newcommand{\tilegrid}[3]{% xoff, title, order macro
\begin{scope}[xshift=#1 cm]
  \node at (3,1.0) {#2};
  \foreach \r in {0,...,5}{\foreach \c in {0,...,5}{
    \pgfmathtruncatemacro{\pid}{#3}
    \ifnum\pid<9 \fill[sysblue!35] (\c,-\r) rectangle ++(1,-1);\fi
    \node at (\c+0.5,-\r-0.5) {\pid};}}
  \draw[step=1] (0,0) grid (6,-6);
\end{scope}}
\tilegrid{0}{row-major order}{\r*6+\c}
% grouped order with GROUP_M=3: pid = group*(3*6) + c*3 + (r mod 3)
\tilegrid{4.5}{grouped order ($G=3$)}{int(\r/3)*18+\c*3+mod(\r,3)}
\node[align=left,anchor=west] at (16.2,-2.0) {shaded: the first 9 programs\\running concurrently};
\node[align=left,anchor=west] at (16.2,-3.6) {row-major: 2 panels of $A$, 6 of $B$\\grouped: 3 panels of $A$, 3 of $B$};
\end{tikzpicture}
\caption{Launch order of output tiles. If 9 programs run at once, grouped ordering
touches 6 input panels instead of 8. For large matrices the saving in L2 misses is
substantial.}
\label{fig:triton:swizzle}
\end{figure}

\begin{lstlisting}[caption={Grouped, tiled matmul. \ident{tl.dot} compiles to tensor-core
MMA instructions on GPUs.},label={lst:triton:matmul}]
@triton.jit
def _matmul_kernel(a_ptr, b_ptr, c_ptr, M, N, K,
                   stride_am, stride_ak, stride_bk, stride_bn, stride_cm, stride_cn,
                   BM: tl.constexpr, BN: tl.constexpr, BK: tl.constexpr,
                   GROUP_M: tl.constexpr):
    pid = tl.program_id(0)
    num_pid_m, num_pid_n = tl.cdiv(M, BM), tl.cdiv(N, BN)
    num_pid_in_group = GROUP_M * num_pid_n
    first_pid_m = (pid // num_pid_in_group) * GROUP_M
    group_size_m = min(num_pid_m - first_pid_m, GROUP_M)
    pid_m = first_pid_m + ((pid % num_pid_in_group) % group_size_m)
    pid_n = (pid % num_pid_in_group) // group_size_m
    rm = pid_m * BM + tl.arange(0, BM)
    rn = pid_n * BN + tl.arange(0, BN)
    rk = tl.arange(0, BK)
    a_ptrs = a_ptr + rm[:, None] * stride_am + rk[None, :] * stride_ak
    b_ptrs = b_ptr + rk[:, None] * stride_bk + rn[None, :] * stride_bn
    acc = tl.zeros((BM, BN), dtype=tl.float32)
    for k0 in range(0, K, BK):
        a = tl.load(a_ptrs, mask=(rm[:, None] < M) & (rk[None, :] + k0 < K), other=0.0)
        b = tl.load(b_ptrs, mask=(rk[:, None] + k0 < K) & (rn[None, :] < N), other=0.0)
        acc += tl.dot(a, b)
        a_ptrs += BK * stride_ak
        b_ptrs += BK * stride_bk
    c_ptrs = c_ptr + rm[:, None] * stride_cm + rn[None, :] * stride_cn
    tl.store(c_ptrs, acc, mask=(rm[:, None] < M) & (rn[None, :] < N))
\end{lstlisting}

\subsection{Pipelining and autotuning}
Inside the $K$ loop, a naive kernel would load a slab, wait, compute, and load the next.
Triton's \code{num\_stages} parameter software-pipelines the loop: the compiler issues
asynchronous copies (\code{cp.async} on Ampere, the Tensor Memory Accelerator on Hopper)
for slabs $k+1,\dots,k+s-1$ while the tensor cores work on slab $k$, keeping $s$ buffers
in shared memory (\cref{fig:triton:pipeline}). \code{num\_warps} sets how many warps
cooperate on one program. Good values depend on the GPU, the shapes and the dtype, so
production kernels declare a list of candidate configurations and let
\code{@triton.autotune} benchmark them the first time each \code{key} (e.g.\ $(M,N,K)$)
is seen:

\begin{lstlisting}[caption={Declaring autotuning configurations (GPU only).},label={lst:triton:autotune}]
@triton.autotune(
    configs=[triton.Config({"BM": 128, "BN": 128, "BK": 64, "GROUP_M": 8},
                           num_stages=3, num_warps=8),
             triton.Config({"BM": 64, "BN": 128, "BK": 32, "GROUP_M": 8},
                           num_stages=4, num_warps=4),
             triton.Config({"BM": 128, "BN": 64, "BK": 32, "GROUP_M": 8},
                           num_stages=4, num_warps=4)],
    key=["M", "N", "K"])
@triton.jit
def _matmul_kernel_tuned(...):  # same body as above
    ...
\end{lstlisting}

\begin{figure}[t]
\centering
\begin{tikzpicture}[x=0.55cm,y=0.55cm,font=\scriptsize]
\newcommand{\ldbox}[3]{\draw[fill=fmemory!40] (#1,#2) rectangle ++(2,0.7); \node at (#1+1,#2+0.35) {#3};}
\newcommand{\mmbox}[3]{\draw[fill=fcompute!35] (#1,#2) rectangle ++(1.4,0.7); \node at (#1+0.7,#2+0.35) {#3};}
\node[anchor=east,font=\small] at (-0.3,3.35) {no pipelining};
\foreach \i in {0,1,2}{\pgfmathsetmacro{\x}{\i*3.4}\ldbox{\x}{3}{load \i}\pgfmathsetmacro{\y}{\x+2}\mmbox{\y}{3}{mma \i}}
\node[anchor=east,font=\small] at (-0.3,1.35) {\texttt{num\_stages}=3};
\foreach \i in {0,1,2,3}{\pgfmathsetmacro{\x}{\i*2}\ldbox{\x}{1.4}{load \i}}
\foreach \i in {0,1,2}{\pgfmathsetmacro{\x}{2+\i*2}\mmbox{\x}{0.4}{mma \i}}
\draw[-Latex] (0,-0.2) -- (12.5,-0.2) node[right] {time};
\node[anchor=west,align=left] at (8.6,1.2) {loads of slabs $k{+}1,k{+}2$ in flight\\while tensor cores run slab $k$};
\end{tikzpicture}
\caption{Software pipelining of the $K$ loop. With multiple shared-memory buffers,
asynchronous loads of future slabs overlap tensor-core work on the current one, hiding
memory latency.}
\label{fig:triton:pipeline}
\end{figure}

\begin{napkinmath}[How close to cuBLAS should you get?]
A square bf16 GEMM with $n=4096$ performs $2n^3=1.37\cdot10^{11}$ FLOPs and touches
$3n^2\cdot2$\,B$=100$\,MB, an intensity of $\approx1365$ FLOP/B, which is deep in
compute-bound territory. At 989\,TFLOP/s the floor is 139\,$\mu$s. cuBLAS typically
reaches 70--80\% of peak at this size on Hopper. A tutorial-style Triton kernel built on
\code{tl.dot} reaches a large fraction of that on Ampere. Matching cuBLAS on Hopper needs
warp-specialised kernels with TMA and asynchronous \code{wgmma}, which newer Triton
versions expose. Do not expect to beat the vendor library on plain GEMMs. Write GEMMs
when you need a \emph{fused prologue or epilogue}, as in the next section.
\end{napkinmath}

% =====================================================================
\section{A fused W8A16 GEMM}\label{sec:triton:w8a16}
\Cref{ch:efficient}'s \code{QuantLinear} dequantized the whole weight in PyTorch before
multiplying, which moves \emph{more} bytes than the bf16 baseline. The fix is to load the
int8 weight tile, widen it in registers, and apply the per-output-channel scale once to
the accumulator in the epilogue. The scale factors out of the $K$ sum because it depends
only on the output channel:
$y_{mn}=\sum_k x_{mk}\,s_n q_{nk}=s_n\sum_k x_{mk}q_{nk}$.

\begin{lstlisting}[caption={Weight-only int8 GEMM: only one byte per weight crosses HBM.},
label={lst:triton:w8a16}]
@triton.jit
def _w8a16_kernel(x_ptr, w_ptr, s_ptr, y_ptr, M, N, K,
                  stride_xm, stride_xk, stride_wn, stride_wk, stride_ym, stride_yn,
                  BM: tl.constexpr, BN: tl.constexpr, BK: tl.constexpr):
    pid_m, pid_n = tl.program_id(0), tl.program_id(1)
    rm = pid_m * BM + tl.arange(0, BM)
    rn = pid_n * BN + tl.arange(0, BN)
    rk = tl.arange(0, BK)
    acc = tl.zeros((BM, BN), dtype=tl.float32)
    for k0 in range(0, K, BK):
        x = tl.load(x_ptr + rm[:, None] * stride_xm + (k0 + rk)[None, :] * stride_xk,
                    mask=(rm[:, None] < M) & ((k0 + rk)[None, :] < K), other=0.0)
        w = tl.load(w_ptr + rn[None, :] * stride_wn + (k0 + rk)[:, None] * stride_wk,
                    mask=(rn[None, :] < N) & ((k0 + rk)[:, None] < K), other=0)
        acc += tl.dot(x.to(tl.float32), w.to(tl.float32))   # widen in registers
    s = tl.load(s_ptr + rn, mask=rn < N, other=0.0)
    acc = acc * s[None, :]                                  # epilogue: apply scale
    tl.store(y_ptr + rm[:, None] * stride_ym + rn[None, :] * stride_yn, acc,
             mask=(rm[:, None] < M) & (rn[None, :] < N))
\end{lstlisting}

Two remarks for real deployments. First, the weight tile is loaded transposed
(\code{w[k, n]} from a row-major \code{(N, K)} array); a production kernel pre-arranges
the layout offline so that loads are contiguous and vectorised. Second, decode has
$M=B$ (a handful of rows), so a $B_M=16$ tile wastes most of each tensor-core
instruction. Specialised kernels (Marlin, Machete) use split-$K$ parallelism so that
enough programs exist to keep every SM streaming weights. For group-wise int4, the scale
\emph{does} depend on $k$, so it is applied per group inside the loop.

% =====================================================================
\section{FlashAttention forward in Triton}\label{sec:triton:flash}
\Cref{ch:attention} derived FlashAttention: process keys and values in blocks, keep a
running row maximum $m$ and normaliser $\ell$, and rescale the partial output whenever
the maximum increases. The $T\times T$ score matrix then never touches HBM. In Triton
the algorithm is short, because tiles are first-class (\cref{lst:triton:flash}).

\begin{proposition}[Online softmax]\label{prop:triton:online}
Process the scores $s_1,\dots,s_S$ of one query row in blocks. Initialise $m=-\infty$,
$\ell=0$, $\vect a=\vect 0$. For each block $\mathcal B$ with values $\vect v_j$ set
$m'=\max(m,\max_{j\in\mathcal B}s_j)$, then
$\ell\leftarrow e^{m-m'}\ell+\sum_{j\in\mathcal B}e^{s_j-m'}$ and
$\vect a\leftarrow e^{m-m'}\vect a+\sum_{j\in\mathcal B}e^{s_j-m'}\vect v_j$, and finally
$m\leftarrow m'$. After the last block, $\vect a/\ell=\sum_j\softmax(\vect s)_j\,\vect v_j$.
\end{proposition}
\begin{proof}
By induction, after processing a set $\mathcal P$ of blocks the invariants
$\ell=\sum_{j\in\mathcal P}e^{s_j-m}$ and $\vect a=\sum_{j\in\mathcal P}e^{s_j-m}\vect v_j$
hold with $m=\max_{j\in\mathcal P}s_j$. Multiplying by $e^{m-m'}$ re-expresses the old
sums relative to the new maximum, and adding the new block's terms extends them to
$\mathcal P\cup\mathcal B$. At the end, $\vect a/\ell=\sum_je^{s_j-m}\vect v_j/\sum_je^{s_j-m}$,
and the common factor $e^{-m}$ cancels. Every exponent is $\le 0$, so nothing
overflows.
\end{proof}

\begin{lstlisting}[caption={FlashAttention forward kernel: one program per (query block,
batch$\times$head).},label={lst:triton:flash}]
@triton.jit
def _flash_fwd_kernel(q_ptr, k_ptr, v_ptr, o_ptr, T, S, sm_scale, stride_b, stride_t,
                      CAUSAL: tl.constexpr, BLOCK_M: tl.constexpr,
                      BLOCK_N: tl.constexpr, D: tl.constexpr):
    pid_m, bh = tl.program_id(0), tl.program_id(1)
    offs_m = pid_m * BLOCK_M + tl.arange(0, BLOCK_M)
    offs_d = tl.arange(0, D)
    base = bh * stride_b
    q = tl.load(q_ptr + base + offs_m[:, None] * stride_t + offs_d[None, :],
                mask=offs_m[:, None] < T, other=0.0).to(tl.float32)
    m_i = tl.full((BLOCK_M,), -float("inf"), dtype=tl.float32)
    l_i = tl.zeros((BLOCK_M,), dtype=tl.float32)
    acc = tl.zeros((BLOCK_M, D), dtype=tl.float32)
    shift = S - T
    hi = S
    if CAUSAL:                                  # skip blocks entirely in the future
        hi = tl.minimum(S, (pid_m + 1) * BLOCK_M + shift)
    for n0 in range(0, hi, BLOCK_N):
        offs_n = n0 + tl.arange(0, BLOCK_N)
        k = tl.load(k_ptr + base + offs_n[:, None] * stride_t + offs_d[None, :],
                    mask=offs_n[:, None] < S, other=0.0).to(tl.float32)
        v = tl.load(v_ptr + base + offs_n[:, None] * stride_t + offs_d[None, :],
                    mask=offs_n[:, None] < S, other=0.0).to(tl.float32)
        s = tl.dot(q, tl.trans(k)) * sm_scale
        valid = offs_n[None, :] < S
        if CAUSAL:
            valid = valid & (offs_n[None, :] <= offs_m[:, None] + shift)
        s = tl.where(valid, s, -float("inf"))
        m_new = tl.maximum(m_i, tl.max(s, axis=1))
        p = tl.exp(s - m_new[:, None])
        alpha = tl.exp(m_i - m_new)             # rescale old statistics
        l_i = l_i * alpha + tl.sum(p, axis=1)
        acc = acc * alpha[:, None] + tl.dot(p, v)
        m_i = m_new
    tl.store(o_ptr + base + offs_m[:, None] * stride_t + offs_d[None, :],
             acc / l_i[:, None], mask=offs_m[:, None] < T)
\end{lstlisting}

\begin{figure}[t]
\centering
\begin{tikzpicture}[x=0.5cm,y=0.5cm,font=\scriptsize]
% score matrix grid of blocks 6x6
\foreach \r in {0,...,5}{\foreach \c in {0,...,5}{
  \ifnum\c>\r \fill[gray!12] (\c,-\r) rectangle ++(1,-1);
  \else\ifnum\c=\r \fill[napkin!45] (\c,-\r) rectangle ++(1,-1);
  \else \fill[sysblue!30] (\c,-\r) rectangle ++(1,-1);\fi\fi}}
\draw[step=1] (0,0) grid (6,-6);
\draw[very thick,warnred] (0,-3) rectangle (6,-4);
\draw[-Latex,thick,warnred] (0,-4.5) -- (4,-4.5) node[midway,below] {loop over key blocks};
\node at (3,0.7) {keys / values $\rightarrow$};
\node[rotate=90] at (-0.8,-3) {query blocks (\texttt{pid\_m})};
\node[anchor=west,align=left] at (7.0,-0.6) {\tikz\fill[sysblue!30] (0,0) rectangle (0.3,0.3); full block: no mask};
\node[anchor=west,align=left] at (7.0,-1.6) {\tikz\fill[napkin!45] (0,0) rectangle (0.3,0.3); diagonal block: causal mask};
\node[anchor=west,align=left] at (7.0,-2.6) {\tikz\fill[gray!12] (0,0) rectangle (0.3,0.3); future block: skipped (\texttt{hi})};
\node[anchor=west,align=left] at (7.0,-4.0) {program for query block 3 (red) keeps\\$m_i$, $\ell_i$, \texttt{acc} on chip; writes\\only its $B_M\times d$ output tile};
\node[anchor=west,align=left] at (7.0,-5.6) {grid $=(\lceil T/B_M\rceil,\;B\cdot H)$};
\end{tikzpicture}
\caption{Work decomposition of the FlashAttention forward kernel. Each program owns one
query block of one head and streams the key/value blocks to its left. With causal
masking, roughly half the blocks are skipped and only diagonal blocks need masks.}
\label{fig:triton:flash}
\end{figure}

\begin{napkinmath}[Why tiling attention matters]
For $T=8192$ and 32 heads, a materialised fp16 score matrix takes
$8192^2\cdot32\cdot2$\,B$=4.3$\,GB \emph{per layer}. It must be written by one kernel and
read by the next (twice more in the backward pass). The tiled kernel's HBM traffic is
only $Q$, $K$, $V$ and $O$, i.e.\ $4\cdot8192\cdot32\cdot128\cdot2$\,B$=268$\,MB, plus
re-reads of $K$ and $V$ by each query block that hit in L2. Attention changes from
memory-bound to compute-bound, and its memory footprint from $O(T^2)$ to $O(T)$.
\end{napkinmath}

\paragraph{Backward pass.} FlashAttention's backward recomputes the probabilities block by
block from $Q$, $K$ and the saved per-row log-sum-exp $m+\log\ell$ (one float per row),
instead of storing $P$. It trades a second $QK^\top$ for $O(T^2)$ memory. The
gradients with respect to $Q$, $K$ and $V$ need a separate pass (or atomics), because
$dK$ and $dV$ accumulate over query blocks.

\paragraph{What later versions add.} FlashAttention-2 (Dao, 2023) swaps the loop order to
parallelise over the sequence dimension even at batch 1, reduces non-matmul FLOPs (the
rescaling), and partitions work between warps to avoid shared-memory communication.
FlashAttention-3 (Shah et al., 2024) targets Hopper: warp specialisation (producer warps
issue TMA loads, consumer warps issue asynchronous \code{wgmma}), ping-pong scheduling
that overlaps the softmax of one tile with the matmul of another, and FP8 support. For
decode, where $T_q=1$, split-KV variants (``flash-decoding'') parallelise over the key
dimension and reduce partial results, because a single query block cannot fill the GPU.

\begin{pitfall}[Our teaching kernel's limits]
\code{flash\_attention} in \code{llmsys} requires $T=S$ (prefill), a power-of-two head
dimension of at least 16, and contiguous \code{(B$\cdot$H, T, D)} layouts. It
accumulates in fp32 but loads whatever dtype it is given. It does not handle GQA head
sharing, paged KV, sliding windows or the backward pass. Use
\code{F.scaled\_dot\_product\_attention} or the FlashAttention library in production.
\end{pitfall}

% =====================================================================
\section{A paged-attention decode kernel}\label{sec:triton:paged}
\Cref{ch:serving} stores the KV cache in fixed-size physical blocks and gives each
sequence a \emph{block table} that maps logical block $j$ to a physical block id. This
eliminates fragmentation, but it means that attention can no longer assume contiguous
keys. The kernel must chase the indirection itself. Decode is the simplest case: one
new query per sequence. \Cref{lst:triton:paged} assigns one program per (sequence,
head), walks the sequence's logical blocks, loads the physical block id from the table,
and applies the online softmax of \cref{prop:triton:online} block by block
(\cref{fig:triton:paged}).

\begin{figure}[t]
\centering
\begin{tikzpicture}[font=\scriptsize,x=0.62cm,y=0.62cm]
% block table
\node[font=\small] at (1.5,1.6) {block table (seq $b$)};
\foreach \j/\p in {0/7,1/2,2/11,3/5}{\draw[fill=sysblue!15] (\j,0) rectangle (\j+1,-0.8); \node at (\j+0.5,-0.4) {\p};
  \node[text=srcgray] at (\j+0.5,0.3) {$j{=}\j$};}
% physical pool
\node[font=\small] at (9.5,1.6) {physical KV blocks (pool)};
\foreach \p in {0,...,11}{
  \pgfmathtruncatemacro{\r}{\p/6}\pgfmathtruncatemacro{\c}{mod(\p,6)}
  \draw[fill=gray!8] (6.4+\c*1.1,-\r*1.3) rectangle ++(1,-1); \node[text=srcgray] at (6.9+\c*1.1,-\r*1.3-0.5) {\p};}
\foreach \p/\col in {7/keygreen,2/keygreen,11/keygreen,5/napkin}{
  \pgfmathtruncatemacro{\r}{\p/6}\pgfmathtruncatemacro{\c}{mod(\p,6)}
  \draw[fill=\col!40] (6.4+\c*1.1,-\r*1.3) rectangle ++(1,-1); \node at (6.9+\c*1.1,-\r*1.3-0.5) {\p};}
\draw[-Latex,thick,fcomm] (0.5,-0.8) to[out=-90,in=-90,looseness=0.6] (8.0,-2.3);
\draw[-Latex,thick,fcomm] (1.5,0) to[out=90,in=90,looseness=0.5] (9.1,0);
\draw[-Latex,thick,fcomm] (2.5,-0.8) to[out=-90,in=-90,looseness=0.5] (12.4,-2.3);
\draw[-Latex,thick,fcomm] (3.5,0) to[out=90,in=90,looseness=0.4] (12.4,0);
\node[anchor=west,align=left] at (0,-4.1) {last block (orange) is partially filled: \texttt{seq\_len} masks its tail.\\
program $(b,h)$ streams blocks 7, 2, 11, 5 with online softmax; output is one $d$-vector.};
\end{tikzpicture}
\caption{Paged decode attention. The block table translates logical KV blocks into
scattered physical blocks. The kernel performs the gather on the fly, so no contiguous
copy of the cache is ever made.}
\label{fig:triton:paged}
\end{figure}

\begin{lstlisting}[caption={Decode attention over a paged KV cache (one query token per
sequence).},label={lst:triton:paged}]
@triton.jit
def _paged_decode_kernel(q_ptr, k_cache_ptr, v_cache_ptr, block_table_ptr, seq_lens_ptr,
                         o_ptr, sm_scale, max_blocks, H: tl.constexpr,
                         BLOCK_SIZE: tl.constexpr, D: tl.constexpr):
    b, h = tl.program_id(0), tl.program_id(1)
    offs_d, offs_t = tl.arange(0, D), tl.arange(0, BLOCK_SIZE)
    q = tl.load(q_ptr + (b * H + h) * D + offs_d).to(tl.float32)
    seq_len = tl.load(seq_lens_ptr + b)
    m_i, l_i = -float("inf"), 0.0
    acc = tl.zeros((D,), dtype=tl.float32)
    for j in range(0, tl.cdiv(seq_len, BLOCK_SIZE)):
        phys = tl.load(block_table_ptr + b * max_blocks + j)     # indirection
        base = (phys * BLOCK_SIZE + offs_t[:, None]) * H * D + h * D + offs_d[None, :]
        valid = (j * BLOCK_SIZE + offs_t) < seq_len
        k = tl.load(k_cache_ptr + base, mask=valid[:, None], other=0.0).to(tl.float32)
        v = tl.load(v_cache_ptr + base, mask=valid[:, None], other=0.0).to(tl.float32)
        s = tl.where(valid, tl.sum(k * q[None, :], axis=1) * sm_scale, -float("inf"))
        m_new = tl.maximum(m_i, tl.max(s, axis=0))
        p = tl.exp(s - m_new)
        alpha = tl.exp(m_i - m_new)
        l_i = l_i * alpha + tl.sum(p, axis=0)
        acc = acc * alpha + tl.sum(p[:, None] * v, axis=0)
        m_i = m_new
    tl.store(o_ptr + (b * H + h) * D + offs_d, acc / l_i)
\end{lstlisting}

Two design notes. First, with a single query the score computation is a matrix--vector
product, so the kernel uses \code{tl.sum} over an elementwise product rather than
\code{tl.dot}. Decode attention is purely memory-bound: every KV byte is read once for
about two FLOPs. The block size therefore matters for \emph{memory} efficiency.
16--32 tokens $\times$ \code{H\_kv} $\times$ \code{D} per block gives contiguous loads
of several KB, enough to approach full bandwidth. Second, the grid has only $B\cdot H$
programs. With a batch of 4 and 8 KV heads, that is 32 programs on 132 SMs, so most of
the GPU idles on long sequences. \emph{Split-KV} (flash-decoding) adds a third grid
axis over chunks of the sequence. Each program returns a partial $(m,\ell,\vect a)$
triple, and a small second kernel merges them with the same rescaling rule. Production
kernels also share the K/V load across the query heads of a GQA group, which cuts KV
traffic by the group size.

\begin{systemslens}[Kernels are the engine's contract]
A serving engine's memory manager (\cref{ch:serving}) and its attention kernel must agree
on everything: block size, KV layout (token-major or head-major), dtype and scale
format for quantized KV, and how the block table and sequence lengths are passed. This is
why engines ship their own attention backends, and why a change such as ``store KV in
FP8'' or ``share blocks across sequences for prefix caching'' is as much a kernel change
as a scheduler change. When you prototype a new cache policy, write the PyTorch
reference (\code{reference.paged\_decode\_attention}) first. It defines the semantics
that any fast kernel must reproduce.
\end{systemslens}

% =====================================================================
\section{Numerics inside kernels}\label{sec:triton:numerics}
Hand-written kernels own their numerics, and the defaults that PyTorch chose for you no
longer apply.
\begin{itemize}
  \item \textbf{Accumulate in fp32.} Tensor cores take bf16/fp16/fp8 inputs but accumulate in
        fp32. Keep \code{acc} in \code{tl.float32} and round once when storing. A bf16
        accumulator over $K=4096$ products loses most of the low-order bits: bf16 has an
        8-bit significand, so once the running sum is 256 times larger than a term, the
        term vanishes.
  \item \textbf{Upcast before reductions.} Our softmax and norm kernels convert loaded tiles
        with \code{.to(tl.float32)} before \code{tl.max}/\code{tl.sum}; fp16 sums of
        squares overflow at 65504.
  \item \textbf{Subtract the max.} \code{exp} of fp32 overflows above $\approx88.7$.
        Every softmax, online or not, must work with $s-m$.
  \item \textbf{Mind the $-\infty$ arithmetic.} A row whose scores are all masked has
        $m=-\infty$, and $e^{m-m'}=e^{-\infty-(-\infty)}$ is NaN. Engines guard this case
        (padding rows in batched decode, fully masked sliding windows), for example by
        clamping $m$ to a finite floor.
  \item \textbf{Determinism.} Floating-point addition is not associative. Split-K GEMMs,
        split-KV attention and atomic accumulations (\code{tl.atomic\_add} for weight
        gradients) sum in an order that depends on scheduling, so results can differ
        bit-for-bit between runs. For RL training this matters: rollouts generated by an
        inference kernel and log-probabilities recomputed by a training kernel can
        disagree slightly even for identical weights (\cref{ch:rlsys}). Deterministic
        modes use fixed reduction trees at some cost in speed.
\end{itemize}

\begin{pitfall}[Testing tolerances]
Compare against an fp32 (or fp64) reference, never against another low-precision kernel,
and scale tolerances with the reduction length. Our tests use $10^{-5}$ for fp32
interpreter runs. The GPU bf16 test allows $2\times10^{-2}$ relative error. A kernel
that passes at $K=64$ and fails at $K=8192$ usually accumulates in the wrong dtype.
\end{pitfall}

\subsection{What to look at in a profile}
\begin{center}\small
\begin{tabularx}{\linewidth}{l>{\raggedright\arraybackslash}X>{\raggedright\arraybackslash}X}
\toprule
kernel type & bound by & first metrics to check (Nsight Compute) \\
\midrule
elementwise, softmax, norms & HBM bandwidth & DRAM throughput (\% of peak); global load efficiency (coalescing) \\
GEMM, prefill attention & tensor cores & tensor pipe utilisation; SMEM bank conflicts; occupancy and register spills \\
decode attention, W4A16 GEMV & HBM bandwidth + parallelism & DRAM throughput; number of programs vs.\ SMs (grid too small?) \\
anything tiny & launch latency & Nsight Systems timeline gaps; use CUDA graphs or fuse \\
\bottomrule
\end{tabularx}
\end{center}

% =====================================================================
\section{Under the hood: from Python to SASS}\label{sec:triton:compiler}
\code{@triton.jit} does not interpret your Python on the GPU. On the first launch with
a new signature it parses the function's AST and lowers it through a stack of MLIR
dialects (\cref{fig:triton:compiler}). \emph{Triton IR} (TTIR) is a
hardware-independent, block-level SSA form. \emph{TritonGPU IR} (TTGIR) attaches
\emph{layouts} to every tensor: which thread owns which elements, whether a tile lives
in registers or in shared memory, and which MMA instruction shape a \code{tl.dot} uses.
Most of Triton's optimisations happen at this level: coalescing, pipelining,
layout conversions, and the insertion of barriers and asynchronous copies. Next come
LLVM IR, then PTX, and finally the vendor assembler produces a cubin (SASS). Every stage
can be inspected:

\begin{lstlisting}[caption={Dumping the intermediate representations of a compiled kernel (GPU).},
label={lst:triton:asm}]
compiled = _softmax_kernel[(n_rows,)](x, out, n_cols, x.stride(0), out.stride(0),
                                      BLOCK=1024)
print(compiled.asm.keys())          # e.g. 'ttir', 'ttgir', 'llir', 'ptx', 'cubin'
print(compiled.asm["ttgir"])        # look for #blocked / #shared / #mma layouts
print(compiled.n_regs, compiled.n_spills)   # register pressure at a glance
\end{lstlisting}

\begin{figure}[t]
\centering
\begin{tikzpicture}[font=\scriptsize,
  st/.style={draw,rounded corners=2pt,fill=#1,minimum height=0.85cm,minimum width=1.55cm,align=center},
  arr/.style={-Latex,thick}]
\node[st=gray!10] (py) {Python\\AST};
\node[st=sysblue!15,right=0.35cm of py] (ttir) {TTIR\\(blocks)};
\node[st=sysblue!30,right=0.35cm of ttir] (ttgir) {TTGIR\\(layouts)};
\node[st=keygreen!20,right=0.35cm of ttgir] (llir) {LLVM IR};
\node[st=keygreen!35,right=0.35cm of llir] (ptx) {PTX};
\node[st=napkin!35,right=0.35cm of ptx] (sass) {cubin\\(SASS)};
\foreach \a/\b in {py/ttir,ttir/ttgir,ttgir/llir,llir/ptx,ptx/sass}{\draw[arr] (\a)--(\b);}
\node[below=0.25cm of ttgir,align=center,text=srcgray] {coalescing, pipelining,\\MMA selection, SMEM allocation};
\node[below=0.25cm of sass,align=center,text=srcgray] {\texttt{ptxas}};
\node[above=0.2cm of ttir,align=center,text=srcgray] {cached per (dtypes, constexprs)};
\end{tikzpicture}
\caption{Triton's compilation pipeline. The optimisations that make block-level code fast
happen when tensors acquire GPU layouts in TTGIR.}
\label{fig:triton:compiler}
\end{figure}

The same applies to code you did not write. \code{torch.compile} with the Inductor
backend emits Triton kernels for fused pointwise and reduction regions (and sometimes
GEMMs, via templates). Set \code{TORCH\_LOGS=output\_code} to print them, or set
\code{TORCH\_COMPILE\_DEBUG=1} to dump them to disk. Reading those kernels is a quick way
to learn what was fused, which regions fell back to library calls, and whether a
\code{view} or \code{.contiguous()} in your model code broke a fusion.

\begin{systemslens}[When the compiler needs help]
Kernels with very large tiles (big \code{BLOCK} for long rows, large GEMM tiles) are
limited by registers. If \code{n\_spills} is non-zero, shrink the tile, lower
\code{num\_stages}, or raise \code{num\_warps} so that each thread holds fewer elements.
Excessive \code{constexpr} specialisation (e.g.\ passing the sequence length as
\code{constexpr}) recompiles for every new value: keep shapes as runtime arguments
unless the compiler genuinely needs them.
\end{systemslens}

% =====================================================================
\section{Triton and CUDA: a translation table}
\begin{center}\small
\begin{tabularx}{\linewidth}{lX}
\toprule
CUDA concept & Triton equivalent / who handles it \\
\midrule
thread block & program instance (\code{tl.program\_id}) \\
grid dims & launch grid tuple \code{kernel[grid](...)} \\
threads per block & \code{num\_warps} $\times$ 32 (compiler distributes the tile) \\
shared memory, \code{\_\_syncthreads} & implicit; inserted by the compiler (layouts, swizzling) \\
coalescing, vector loads & inferred from contiguous \code{tl.arange} offsets \\
\code{mma}/\code{wgmma} fragments & \code{tl.dot} on tiles $\ge16$ in each dimension \\
\code{cp.async}, TMA, multi-buffering & \code{num\_stages}; tensor descriptors in recent versions \\
warp shuffles for reductions & \code{tl.sum}, \code{tl.max} along an axis \\
atomics & \code{tl.atomic\_add}, \code{tl.atomic\_max}, \dots \\
templates / compile-time params & \code{tl.constexpr} arguments (re-JIT per value) \\
\bottomrule
\end{tabularx}
\end{center}
When does a project drop below Triton? When it needs explicit control over warp roles,
register allocation or inter-SM cooperation (thread-block clusters, distributed shared
memory), or the very last 10--20\% on a flagship GEMM or attention kernel. Then use CUDA
C++ with CUTLASS/CuTe, or tile DSLs such as ThunderKittens.

% =====================================================================
\section{Measuring kernels}
\begin{itemize}
  \item \textbf{Wall-clock}: \code{triton.testing.do\_bench(fn)} handles warm-up, L2
        flushing and CUDA-event timing. Never time a GPU kernel with
        \code{time.time()} without synchronising.
  \item \textbf{Achieved rates}: convert time into GB/s for memory-bound kernels and TFLOP/s
        for compute-bound ones, and report the percentage of the roofline (the lab's
        \code{bench.py} does this). A 2$\times$ speedup that reaches only 15\% of peak
        bandwidth still has headroom.
  \item \textbf{Timelines}: Nsight Systems (\code{nsys profile python script.py}) shows gaps
        between kernels: launch overhead, CPU stalls, synchronisations. For decode, these
        gaps are often the real problem, and the fix is CUDA graphs (\cref{ch:inference}),
        not a faster kernel.
  \item \textbf{Inside one kernel}: Nsight Compute (\code{ncu --set full}) reports DRAM
        throughput, SM and tensor-pipe utilisation, occupancy, bank conflicts and stall
        reasons, and plots the kernel on a roofline.
\end{itemize}

\subsection{Using a Triton kernel from PyTorch}
A launcher like \code{softmax\_rows} is already callable on CUDA tensors. To make it a
first-class operator that \code{torch.compile} can trace through and autograd can
differentiate, register it as a custom op (\code{torch.library.custom\_op}) with a fake
(meta) implementation for shape inference, and give it a backward formula, either with
\code{register\_autograd} or by wrapping it in a \code{torch.autograd.Function}. Recent
PyTorch versions can also capture user-defined Triton kernels directly inside compiled
graphs.

% =====================================================================
\section{Summary}
\begin{itemize}
  \item Triton programs operate on tiles. You choose the decomposition (grid, block sizes,
        loop order), and the compiler handles threads, shared memory and tensor cores.
  \item Masks with neutral padding values (\code{0}, $-\infty$) make ragged edges and
        reductions correct.
  \item Fusion turns memory-bound chains into single passes. Tiling changes asymptotic
        traffic (GEMM, FlashAttention). Program ordering and pipelining decide how close a
        tiled kernel gets to the roofline.
  \item Quantized GEMMs and attention are where custom kernels pay off most in LLM
        inference, because their data formats and fusions are not library defaults.
  \item Develop with the interpreter and PyTorch references, then measure against the
        roofline with \code{do\_bench}, Nsight Systems and Nsight Compute.
\end{itemize}

% =====================================================================
\section{Exercises}
\begin{exercise}[grid arithmetic]
For \code{vector\_add} with $n=10^6$ and \code{BLOCK}=1024, how many programs are
launched, and how many lanes are masked off? Why must \code{BLOCK} be a power of two?
\end{exercise}
\begin{exercise}[padding values]
In \cref{lst:triton:softmax}, what goes wrong if masked lanes are loaded with
\code{other=0.0}? Construct a row of length 3 with \code{BLOCK}=4 where the result is
incorrect.
\end{exercise}
\begin{exercise}[long rows]
Rewrite the softmax kernel for rows longer than one block, using a loop with online
max/sum rescaling followed by a second normalisation pass. How many HBM passes does your
kernel make? Can you get away with one write?
\end{exercise}
\begin{exercise}[tile intensity]
Compute $I_{\text{tile}}$ for $(B_M,B_N)\in\{(64,64),(128,128),(128,256)\}$ in bf16 and
in fp8. Which of them can saturate an H100's tensor cores purely from HBM bandwidth, and
why does L2 reuse change the answer?
\end{exercise}
\begin{exercise}[swizzling]
Generalise \cref{fig:triton:swizzle}: for an $R\times C$ grid of tiles, $P$ concurrent
programs and group size $G$, count the distinct $A$ and $B$ panels touched by the first
$P$ programs, and find the $G$ that minimises the total.
\end{exercise}
\begin{exercise}[shared memory budget]
With \code{num\_stages}$=s$ and bf16 tiles $B_M\times B_K$ and $B_K\times B_N$, how
much shared memory does the pipelined GEMM need? With 228\,KB per SM on an H100, what is
the largest $s$ for $128\times128\times64$ tiles, and how many such programs fit per SM?
\end{exercise}
\begin{exercise}[int4 epilogue]
Modify the W8A16 kernel for group-wise int4 with packed weights: where do you unpack,
and where must the scale and zero point be applied when the group size equals $B_K$?
\end{exercise}
\begin{exercise}[FlashAttention memory]
Count the HBM bytes read and written by the forward kernel for $B=1$, $H=32$,
$T=S=8192$, $d=128$, bf16, assuming no L2 hits, as a function of $B_M$. Compare with
naive attention.
\end{exercise}
\begin{exercise}[decode attention]
Explain why the kernel of \cref{lst:triton:flash} uses only $B\cdot H$ programs when
$T=1$, and sketch a split-KV variant: what extra state must each partial result carry so
the reduction is exact?
\end{exercise}

\begin{exercise}[SwiGLU in the epilogue]
Sketch how to fuse SiLU-and-multiply into the epilogue of a GEMM that computes
\code{gate} and \code{up} together from a concatenated weight $[W_g;W_u]$. Which tile
shape constraint does this impose, and how many HBM bytes per token does it save for
Qwen3-8B?
\end{exercise}
\begin{exercise}[RMSNorm backward]
Derive \cref{eq:triton:rmsbwd}, starting from
$\partial\hat x_i/\partial x_j=r\big(\delta_{ij}-\hat x_i\hat x_j/n\big)$. Why is it
safe, and cheaper, to recompute $r$ in the backward kernel rather than save it?
\end{exercise}
\begin{exercise}[paged decode]
For Qwen3-8B (8 KV heads, $d=128$, bf16) and a block size of 16 tokens, how many bytes
does one program load per block? With a batch of 2 and 32k-token contexts, how many
programs does \cref{lst:triton:paged} launch, and how many split-KV chunks per sequence
would you use on a 132-SM GPU?
\end{exercise}
% =====================================================================
\section{Lab A: GPU kernels with Triton}\label{sec:triton:lab}
\begin{labbox}{A}{GPU kernels with Triton}
\labdir{labA_triton}
\labsection{Goals}
Write five Triton kernels (vector add, fused softmax, RMSNorm, grouped tiled matmul,
FlashAttention forward) and a PyTorch online-softmax attention reference. Verify them on
any machine with the interpreter, then measure them against the roofline on a GPU.

\labsection{Setup}
\code{cd code; LAB\_IMPL=starter python -m pytest -q labs/labA\_triton}. Without CUDA,
the test file sets \code{TRITON\_INTERPRET=1} automatically (use small shapes; it is
slow). \code{\_w8a16\_kernel} in \code{starter.py} is complete: read it as a worked
example before you start.

\labsection{Tasks}
\begin{enumerate}
  \item \code{\_add\_kernel} and \code{vector\_add}: the grid, offsets and masks.
  \item \code{\_softmax\_kernel} and \code{softmax\_rows}: one program per row; choose
        padding values that are neutral for max and sum.
  \item \code{\_rmsnorm\_kernel} and \code{rmsnorm}: mean of squares in fp32; support
        arbitrary leading dimensions by flattening in the launcher.
  \item \code{\_matmul\_kernel} and \code{matmul}: tiles, a $K$ loop with masks, fp32
        accumulation with \code{tl.dot}, then grouped program ordering.
  \item \code{flash\_attention\_tiled\_torch}: online softmax in PyTorch, including
        $T\neq S$ (queries aligned to the end of the keys, as in decode).
  \item \code{\_flash\_fwd\_kernel} and \code{flash\_attention}: port your PyTorch version
        to Triton, with causal block skipping.
  \item \code{arithmetic\_intensity} and \code{roofline\_time}: the napkin math used in
        \code{bench.py}.
\end{enumerate}

\labsection{Measurements (GPU)}
Run \code{bench.py --impl starter --hw <your GPU>}. Report, for each kernel and size,
time, achieved GB/s or TFLOP/s and percentage of peak. Expect vector add and softmax at
70--90\% of bandwidth for large sizes, and the unfused eager softmax several times
slower. For matmul, compare with and without grouped ordering (\code{group\_m=1}) and
for 2--3 tile configurations; then add \code{@triton.autotune}. For FlashAttention,
compare against \code{F.scaled\_dot\_product\_attention} and plot time against $T$ for
your kernel and for naive attention: one curve should be quadratic, and you should
explain why the other is too, despite linear memory.

\labsection{Deliverables}
Passing tests (interpreter or GPU); a table of achieved rates vs.\ roofline; one
paragraph on the biggest gap you found and its cause, backed by an Nsight Compute
metric.

\labsection{Stretch goals}
\begin{itemize}
  \item Fuse RoPE into the FlashAttention kernel's $Q$ and $K$ loads.
  \item Add GQA: let several query heads (\code{bh} indices) map to one KV head without
        materialising repeated K/V.
  \item Implement a split-KV decode kernel and compare it with the prefill kernel at $T=1$.
  \item Write the group-wise int4 variant of the W8A16 kernel and plug it into
        \code{QuantLinear} from Lab 13. Measure decode tokens/s for Qwen3-0.6B.
\end{itemize}
\end{labbox}
